\begin{frame}{OOB(out-of-bag): 검증 세트 없이도 일반화 성능을}
  \begin{columns}[T]
    \begin{column}{0.46\textwidth}
      \begin{itemize}
        \item 그 37\% = \kb{OOB}(out-of-bag, bag 밖) 샘플 — 해당 트리가
          \textbf{보지 못한} 샘플
        \item 이 트리가 ``보지 못한'' 샘플을 예측하게 하면, 학습을
          마치자마자 \textbf{새 데이터에 대한 예측}을 할 수 있다 —
          \textbf{별도 검증 세트를 떼어놓을 필요가 없음}
      \end{itemize}
      \vskip0.5em
      \begin{block}{실무적 중요성}
        \small \textbf{데이터가 적어 검증 세트를 따로 떼기 아까울 때}에도
        ``이 모델이 과적합이 아닌가''를 \textbf{학습 중 즉시} 확인할 수
        있다.
      \end{block}
    \end{column}
    \begin{column}{0.5\textwidth}
      \centering
      \includegraphics[width=\linewidth]{ch07_2_oob_diagram.png}
      \vskip0.3em
      {\footnotesize 복원추출로 뽑은 부트스트랩 샘플(약 63\%)과 한 번도
      뽑히지 않은 OOB 샘플(약 37\%).}
    \end{column}
  \end{columns}

\note{\begin{itemize}
\item 각 트리는 \textbf{복원추출}된 부트스트랩 샘플로 학습하는데,
      복원추출은 원본 샘플 중 \textbf{약 37\%를 ``한 번도 안 고른다''}
\end{itemize}}
\end{frame}
